\subsection*{寻找扩散时不改变形状的波}
任意温度曲线难追踪，但正弦有一个特殊性质：对空间求两次导数，仍是自身乘一个负常数。若初始形状是 $\sin(kx)$，热方程只要求振幅随时间满足 $A'=-\kappa k^2A$，所以振幅指数衰减，形状保持。

这与线性代数的特征方向相似。求二阶导是一个线性算子，正弦是它的特征模式；热演化在每个模式上只改变系数。对固定零温边界，频率还必须使正弦在两个端点都为零，因此允许的 $k$ 不是任意数。

\subsection*{高频为什么先消失}
频率增为三倍，二阶导倍率增为九倍，因此高频振幅衰减得更快。温度曲线上细小急促的起伏先被抹平，宽缓结构保留更久；“热让分布平滑”由此有了可计算的原因。

动画采用周期边界的一低一高两个模式：蓝色是低频，黄色是高频，白色是总和。先在暂停时预测哪部分会先消失，再推进时间，观察总曲线逐渐接近低频形状。这里的边界选择与正文零温杆示例不同，但模式衰减机制相同。

有限个模式相加可以直接利用线性性；无限级数则需要额外条件，才能交换求和、求导并恢复初值。把每个模式代回方程是必要检查，却还没有自动完成一般解理论。这一点将在傅里叶章节继续讨论。
